\documentclass{article}

\usepackage{Paper}

\fancyhead[L]{EE\_ECO}
\pdftitle{Life Cycle Assessment}
\setlength{\droptitle}{-2.5em}

% === TITLE ===
\title{\textbf{Life cycle assessment of different energy technologies \\ \Large Environmental Analysis \& Ecology}}
\author{Matteo Frongillo, Anastasia Linnik, Yannik Neukom}
\date{}

% === TEXT ===
\begin{document}
\maketitle
\thispagestyle{fancy}

\section{Photovoltaic Park {\small (Matteo)}}
\subsection{Paper 1}
\textbf{Title:} Estimating the environmental footprint of a grid-connected 20 MWp photovoltaic system\\
\textbf{Authors:} Artúr Szilágyi, Gyula Gróf\\
\textbf{Year:} 2020\\
\textbf{Journal/Source:} Solar Energy \parencite{SZILAGYI2020491}

\dots

\subsection{Paper 2}
\textbf{Title:} Life Cycle Assessment of a ground-mounted 1778 kW$_\text{p}$ photovoltaic plant and comparison with traditional energy production systems\\
\textbf{Authors:} Umberto Desideri, Stefania Proietti, Francesco Zepparelli, Paolo Sdringola, Silvia Bini\\
\textbf{Year:} 2012\\
\textbf{Journal/Source:} Applied Energy \parencite{DESIDERI2012930}

\dots

\subsection{Comparison of Paper 1 and Paper 2}

\dots

\newpage
\section{Nuclear Power Plant {\small (Yannik)}}
\subsection{Paper 1}
\textbf{Title:} \\
\textbf{Authors:} \\
\textbf{Year:} \\
\textbf{Journal/Source:} 

\dots

\subsection{Paper 2}
\textbf{Title:} \\
\textbf{Authors:} \\
\textbf{Year:} \\
\textbf{Journal/Source:}

\dots

\subsection{Comparison of Paper 1 and Paper 2}

\dots

\newpage
\section{Coal-fired Power Plant {\small (Anastasia)}}
\subsection{Paper 1}
\textbf{Title:} \\
\textbf{Authors:} \\
\textbf{Year:} \\
\textbf{Journal/Source:} 

\dots

\subsection{Paper 2}
\textbf{Title:} \\
\textbf{Authors:} \\
\textbf{Year:} \\
\textbf{Journal/Source:} 

\dots

\subsection{Comparison of Paper 1 and Paper 2}

\dots
















\newpage
\appendix
\section{LCA Questionnaire}
\subsection{Goal and scope}
\subsubsection{Goal of LCA}
\begin{enumerate}
    \renewcommand{\labelenumi}{\alph{enumi})}
    \item What is the intended application of the LCA study?
    \item What is the purpose of the LCA study?
    \item Who is the intended audience of the LCA report?
    \item Is the LCA study based on a comparative analysis?
\end{enumerate}

\subsubsection{Scope of LCA}
\begin{enumerate}
    \renewcommand{\labelenumi}{\alph{enumi})}
    \setcounter{enumi}{4}
    \item What is the function(s) of the product and which functional unit has been applied?
    \item Plan and describe the system boundaries of the LCAs using a process flow diagram showing the processes, their relationships, and the function of the system.
    \item What data and type of data are needed to calculate inventory results
    and selected impact assessment categories?
    \begin{itemize}[label=--]
        \item Data quality requirements, assumptions and limitations
        \item Data acquisition: Is the data measured, calculated or estimated? How much of the data required is primary data (in \%) and how much data is taken from literature and databases (secondary data)?
        \item Time-reference: When was the data obtained and have there been any major changes since the data collection that might affect the results?
        \item Geographical reference: For what country or region is this data relevant?
        \item Is the secondary data from literature or databases representative for state-of-the-art-technology or for older technology?
        \item Precision: Is the data a precise representation of the system?
        \item Completeness: Are any data missing? How are data gaps filled?
        \item Representativeness, consistency, reproducibility: Is the data representative, consistent and can it be reproduced?
    \end{itemize}
    \item What impact categories are used?
\end{enumerate}

\subsection{Results of Life Cycle Impact Assessment (LCIA)}
\begin{enumerate}
    \renewcommand{\labelenumi}{\alph{enumi})}
    \item State the results of the different impact categories in numbers.
    \item Interpret the two case studies by considering the following:\\
        Results are checked and evaluated to check if they are consistent with the goal and scope definition. This includes: identification of significant issues
        (e.g. inventory elements or impact categories missing) and evaluation/interpretation of results (e.g. comparability)
\end{enumerate}

\newpage
\section{References}
\printbibliography



\end{document}
